\documentclass[white,11pt,oneside]{oritone-manuscript}

\usepackage{array}
\lstset{mathescape=false}

\title{Oritone LaTeX Manual}
\author{Oritone Maintainers}
\date{2026-05-04}

\newcommand{\pkg}[1]{\textsf{#1}}
\newcommand{\cls}[1]{\textsf{#1}}
\newcommand{\opt}[1]{\texttt{#1}}
\newcommand{\cmd}[1]{\texttt{\textbackslash{}#1}}
\newcommand{\envname}[1]{\texttt{#1}}
\newcommand{\filename}[1]{\begingroup\Urlmuskip=0mu plus 2mu\path{#1}\endgroup}
\newcolumntype{L}[1]{>{\raggedright\arraybackslash}p{#1}}

\begin{document}
\maketitle

\begin{abstract}
Oritone is a LaTeX package and class collection for CVs, manuscripts, theses,
and Beamer presentations. It centralizes color roles, font roles, listings,
plots, localized labels, callouts, title rules, and page furniture so the
individual document frontends share one visual language.
\end{abstract}

\keywords{LaTeX\keywordsep document classes\keywordsep Beamer theme\keywordsep
Recursive fonts\keywordsep technical documents}

\tableofcontents

\section{Package Layout}

The public source tree is intentionally small:

\begin{center}
\small
\renewcommand{\arraystretch}{1.15}
\begin{tabular}{L{0.40\linewidth}L{0.49\linewidth}}
\hline
Path & Purpose \\
\hline
\filename{tex/latex/oritone/oritone.sty} &
Shared palette, color modes, semantic color roles, listings, and PGFPlots
integration. \\
\filename{tex/latex/oritone/oritone-fonts.sty} &
Engine-aware font role loading for serif, sans, mono, and math fonts. \\
\filename{tex/latex/oritone/oritone-design.sty} &
Shared document-design helpers: labels, links, page colors, callouts, and
ruled title blocks. \\
\filename{tex/latex/oritone/oritone-cv.cls} &
Compact CV class based on \cls{article}. \\
\filename{tex/latex/oritone/oritone-manuscript.cls} &
Review-friendly manuscript class based on \cls{extarticle}. \\
\filename{tex/latex/oritone/oritone-thesis.cls} &
Long-form thesis class based on \cls{book}. \\
\filename{tex/latex/oritone/beamerthemeOritone.sty} and subthemes &
Beamer theme, color theme, font theme, inner theme, outer theme, and shared
theme options. \\
\filename{examples/} &
Filled example documents and mode wrappers used by the documentation build. \\
\filename{testfiles/} &
\pkg{l3build} regression and smoke tests. \\
\hline
\end{tabular}
\end{center}

\section{Installation}

From a checkout, install the files into the local TeX tree with:

\begin{lstlisting}
l3build install
\end{lstlisting}

For one-off local testing without installing, put the source directory on
\texttt{TEXINPUTS}:

\begin{lstlisting}
TEXINPUTS="$(pwd)/tex/latex/oritone//:" pdflatex examples/manuscript-white.tex
\end{lstlisting}

The package targets LaTeX2e from 2020-10-01 or newer. The examples and tests
are expected to compile with pdfLaTeX, LuaLaTeX, and XeLaTeX, subject to the
font notes in \Cref{sec:fonts}.

\section{Core Package}

The base package is loaded directly when a document only needs the palette,
listing style, or plot style:

\begin{lstlisting}[language={[LaTeX]TeX}]
\usepackage[mode=white,accent=blue,listings,plots]{oritone}
\end{lstlisting}

\subsection{Options}

\begin{center}
\small
\renewcommand{\arraystretch}{1.15}
\begin{tabular}{p{0.24\linewidth}p{0.22\linewidth}p{0.43\linewidth}}
\hline
Option & Values & Effect \\
\hline
\opt{mode} & \opt{light}, \opt{dark}, \opt{white}, \opt{black} &
Selects the active foreground, background, surface, and semantic color roles. \\
\opt{accent} &
\opt{auto}, \opt{orange}, \opt{sky}, \opt{green}, \opt{yellow}, \opt{blue},
\opt{vermillion}, \opt{rose}, \opt{purple}, \opt{red}, \opt{amber},
\opt{cyan} &
Selects the accent hue used for links, rules, headings, and progress bars. \\
\opt{accentshade} &
\opt{normal}, \opt{dark}, \opt{darker}, \opt{light}, \opt{lighter} &
Shades the whole accent family toward the foreground text (\opt{dark}/\opt{darker})
or the page background (\opt{light}/\opt{lighter}) so links, rules, headings, the
footer bar, and progress bars shift together. Useful when the default accent
looks too bright on the light and white modes. \\
\opt{listings} or \opt{julia} & Boolean &
Loads \pkg{listings} and enables the shared Oritone listing styles. \\
\opt{plots} or \opt{pgfplots} & Boolean &
Loads \pkg{pgfplots} and enables the shared Oritone plot cycle and axis
defaults. \\
\hline
\end{tabular}
\end{center}

The active mode and accent can also be changed after package load:

\begin{lstlisting}[language={[LaTeX]TeX}]
\oritonesetup{mode=dark,accent=amber}
\oritonesetup{mode=white,accent=blue,accentshade=dark}
\end{lstlisting}

\section{Fonts}
\label{sec:fonts}

Font loading is centralized in \pkg{oritone-fonts}. The default role profile is:

\begin{center}
\small
\renewcommand{\arraystretch}{1.15}
\begin{tabular}{p{0.2\linewidth}p{0.32\linewidth}p{0.37\linewidth}}
\hline
Role & Default & Used for \\
\hline
Serif & STIX Two Text & Body text in manuscript and thesis documents. \\
Math & STIX Two Math & Math in unicode engines; pdfLaTeX uses the STIX2
package. \\
Sans & Roboto Condensed & CV body text, headings, title pages, headers, footers,
Beamer text, labels, and other display roles. \\
Mono & Recursive Mono & Listings, verbatim-like text, and code-oriented helper
macros. \\
\hline
\end{tabular}
\end{center}

LuaLaTeX and XeLaTeX use \pkg{fontspec} and \pkg{unicode-math}. They first try
installed Roboto, Recursive, and STIX names, then known static font file names,
then fallback families.

pdfLaTeX cannot load OpenType system fonts directly. For Roboto Condensed it
uses the \pkg{roboto} package, and for Recursive Mono the optional external
\pkg{recursive} package when \filename{recursive.sty} is installed. If a support
package is missing, Oritone warns and falls back to PT Sans and Inconsolata so
the document remains compileable.

\begin{lstlisting}
git clone git@github.com:schmaeke/latex-recursive.git
cd latex-recursive
l3build install
\end{lstlisting}

Advanced users can load \pkg{oritone-fonts} directly with role overrides:

\begin{lstlisting}[language={[LaTeX]TeX}]
\usepackage[
  serif=TeX Gyre Termes,
  sans=Recursive Sans,
  mono=Recursive Mono,
  math=STIX Two Math,
  sansscale=0.95,
  monoscale=1.0
]{oritone-fonts}
\end{lstlisting}

\section{Shared Design Roles}

\pkg{oritone-design} is loaded by the classes and theme. It provides:

\begin{itemize}
  \item \cmd{oritoneApplyPageColors} for applying the active foreground and
    background to the page.
  \item \cmd{oritoneSetSemanticLinks} for semantic \pkg{hyperref} colors.
  \item \cmd{oritoneRuledTitle} for title text framed by horizontal rules.
  \item \cmd{oritonedefinecallouts} for \envname{oritonenote},
    \envname{oritoneexample}, and \envname{oritonewarning}.
  \item \cmd{oritonelabel} and \cmd{oritonecurrentlabel} for localized labels
    such as page, abstract, keywords, declaration, and footer text.
\end{itemize}

Most documents should not load \pkg{oritone-design} directly. Use the classes
or the Beamer theme unless you are building another Oritone-compatible frontend.

\section{CV Class}

The CV class is a compact, sans-serif \cls{article} derivative:

\begin{lstlisting}[language={[LaTeX]TeX}]
\documentclass[white,10pt,a4paper,language=english]{oritone-cv}

\name{Ada}
\surname{Sterling}
\shortinfo{\maillink{ada@example.com}\sep\weblink{\faGlobe}{ada.dev}{ada.dev}}
\phrase{I build reliable research software.}

\begin{document}
\cvhead
\cvsection{Experience}
\cvdatedentry{Example Laboratory}{2022--Present}{Built reusable workflows.}
\end{document}
\end{lstlisting}

\subsection{CV Options}

The class accepts \opt{light}, \opt{dark}, \opt{white}, \opt{black}, or
\opt{mode=...}; \opt{accent=...}; \opt{language=...}; ordinary \cls{article}
paper options; and \opt{10pt}, \opt{11pt}, or \opt{12pt}. If no paper or size is
selected, it defaults to A4 and 10 pt.

\subsection{CV Commands}

\begin{center}
\small
\renewcommand{\arraystretch}{1.15}
\begin{tabular}{p{0.35\linewidth}p{0.52\linewidth}}
\hline
Command & Purpose \\
\hline
\cmd{name}, \cmd{surname}, \cmd{shortinfo}, \cmd{phrase} &
Store header metadata for \cmd{cvhead}. \\
\cmd{cvhead} & Render the name, contact line, accent rule, and phrase. \\
\cmd{cvsection} & Schedule a section heading. \\
\cmd{cventry} & Entry with title and body. \\
\cmd{cvdatedentry} & Entry with title, date/meta column, and body. \\
\cmd{cvdatedentryshort} & Compact dated entry for projects, talks, skills, or
short lists. \\
\cmd{cvcustomentry} & Explicit title, right-side meta, and body. \\
\cmd{maillink}, \cmd{phonelink}, \cmd{weblink}, \cmd{inlinelink} &
Contact and inline link helpers. \\
\cmd{sep}, \cmd{lastupdated}, \cmd{highlighted} &
Inline separators and small text helpers. \\
\hline
\end{tabular}
\end{center}

\section{Manuscript Class}

The manuscript class is intended for review drafts, internal reports, preprints,
and short technical papers:

\begin{lstlisting}[language={[LaTeX]TeX}]
\documentclass[white,11pt,oneside,accent=blue]{oritone-manuscript}

\title{A Reusable Scientific Report}
\author[1]{Ada Example}
\affil[1]{Example University}
\date{\today}

\begin{document}
\maketitle

\begin{abstract}
The abstract spans the full text width and has no first-line indent.
\end{abstract}

\keywords{reports\keywordsep LaTeX\keywordsep reusable classes}

\section{Introduction}
Body text starts here.
\end{document}
\end{lstlisting}

\subsection{Manuscript Options}

The class accepts \opt{light}, \opt{dark}, \opt{white}, \opt{black}, or
\opt{mode=...}; \opt{accent=...}; \opt{language=...}; \opt{oneside} or
\opt{twoside}; common paper options; and \cls{extarticle} sizes from
\opt{8pt} through \opt{20pt}. The default is A4, 10 pt, two-sided.

\subsection{Manuscript Features}

\begin{itemize}
  \item Full-width abstract with localized sans-serif heading.
  \item \cmd{keywords} helper with localized label and \cmd{keywordsep}
    separator.
  \item Optional \cmd{printcompileinfo} footer timestamp.
  \item Dense sans-serif headings through \pkg{titlesec}.
  \item \pkg{authblk} author and affiliation support.
  \item \pkg{lineno}, \pkg{cleveref}, \pkg{hyperref}, \pkg{listings}, TikZ, and
    semantic callouts.
\end{itemize}

\section{Thesis Class}

The thesis class is a longer-form \cls{book} derivative with front matter,
metadata helpers, bibliography helpers, declaration page, plots, listings, and
two-sided running heads.

\begin{lstlisting}[language={[LaTeX]TeX}]
\documentclass[english,white,bibtex]{oritone-thesis}

\title{A Thesis Built with Oritone}
\author{Ada Example}
\date{\today}

\universityName{Example University}
\studyProgram{Computational Engineering}
\desiredDegree{Master of Science}
\firstExaminer{Prof. A. Example}{Computational Methods}
\secondExaminer{Prof. B. Sample}{Software Systems}
\thesisSubject{Thesis class demonstration}
\thesisKeywords{LaTeX, thesis, Oritone}
\thesisBibliographyResource{references}

\begin{document}
\frontmatter
\maketitle
\abstractChapter
Abstract text.
\tableofcontents
\mainmatter
\chapter{Introduction}
\printThesisBibliography
\declarationOfAcademicIntegrity
\end{document}
\end{lstlisting}

\subsection{Thesis Options}

The class accepts \opt{light}, \opt{dark}, \opt{white}, \opt{black}, or
\opt{mode=...}; \opt{accent=...}; \opt{language=...}; \opt{oneside} or
\opt{twoside}; common paper options; \opt{10pt}, \opt{11pt}, or \opt{12pt};
\opt{draft} or \opt{final}; \opt{bibtex} or \opt{biblatex}; and
\opt{colorlinks} or \opt{hidelinks}. The default is A4, 12 pt, two-sided,
final, BibTeX, and color links.

\subsection{Thesis Commands}

\begin{center}
\small
\renewcommand{\arraystretch}{1.15}
\begin{tabular}{p{0.36\linewidth}p{0.51\linewidth}}
\hline
Command & Purpose \\
\hline
\cmd{universityName}, \cmd{universityLogo}, \cmd{universityLogoWidth} &
Institution title-page metadata. \\
\cmd{studyProgram}, \cmd{desiredDegree} & Degree metadata. \\
\cmd{firstExaminer}, \cmd{secondExaminer} & Examiner metadata. \\
\cmd{thesisSubject}, \cmd{thesisKeywords} & PDF metadata. \\
\cmd{abstractChapter} & Localized unnumbered abstract chapter. \\
\cmd{thesisBibliographyStyle}, \cmd{thesisBibliographyResource},
\cmd{printThesisBibliography} & BibTeX/BibLaTeX abstraction helpers. \\
\cmd{declarationOfAcademicIntegrity} & Localized declaration page. \\
\cmd{todo}, \cmd{linktt}, \cmd{intd}, \cmd{Assembly}, \cmd{bs} &
Small thesis prose and math helpers. \\
\hline
\end{tabular}
\end{center}

\section{Beamer Theme}

Use the theme with a regular Beamer document:

\begin{lstlisting}[language={[LaTeX]TeX}]
\documentclass[aspectratio=1610]{beamer}
\usetheme[mode=light,accent=blue]{Oritone}

\title{A Complete Oritone Deck}
\author{Ada Example}
\institute{Example Laboratory}
\date{\today}
\biglogo{example-image-a}
\smalllogo{example-image-b}

\begin{document}
\maketitle
\sectionframe{Overview}
\begin{frame}{Motivation}
  \begin{block}{Claim}
    Shared package roles keep slides and documents consistent.
  \end{block}
\end{frame}
\closingframe{Thank you}{example.com/oritone}{Questions}
\end{document}
\end{lstlisting}

\subsection{Beamer Options}

The theme accepts \opt{light}, \opt{dark}, \opt{white}, \opt{black}, or
\opt{mode=...}; \opt{accent=...}; \opt{accentshade=...}; the accent shortcut
options; and \opt{draftlogos} or \opt{showlogoplaceholders}. Missing logo
placeholders are silent unless a placeholder option is enabled.

\subsection{Beamer Commands}

\begin{center}
\small
\renewcommand{\arraystretch}{1.15}
\begin{tabular}{p{0.35\linewidth}p{0.52\linewidth}}
\hline
Command & Purpose \\
\hline
\cmd{titleframe} or \cmd{maketitle} & Render the Oritone title slide. \\
\cmd{sectionframe} & Create a section divider slide. \\
\cmd{standoutframe} & Create a slide without normal header/footer chrome and
with a left accent rule. \\
\cmd{closingframe} & Render a closing slide with title and optional logo. \\
\cmd{slidebackground}, \cmd{noslidebackground} & Set or clear the optional
curved slide background image. \\
\cmd{coauthors}, \cmd{biglogo}, \cmd{smalllogo} & Presentation metadata helpers. \\
\cmd{weblink} & Link helper for bare host names. \\
\hline
\end{tabular}
\end{center}

\section{Listings and Plots}

Classes call \cmd{oritonelistingssetup} and select the shared
\opt{oritone-listing} style automatically. The base style uses Recursive Mono
or the configured mono fallback, a smaller code size, and the document's line
spacing stretch.

\begin{lstlisting}[language=Julia,caption={A compact listing.}]
function residual(values)
  return sum(abs, values)
end
\end{lstlisting}

The thesis class and Beamer examples call \cmd{oritonepgfplotssetup} to install
the shared axis and color-cycle defaults. Users of the base package can enable
the same behavior with:

\begin{lstlisting}[language={[LaTeX]TeX}]
\usepackage[plots]{oritone}
\oritonepgfplotssetup
\end{lstlisting}

The plot cycle follows the Okabe-Ito colorblind-safe palette, mapped to the
closest Oritone hues and ordered blue, vermillion, green, purple, sky, orange,
yellow, and gray. The colors are available individually as \opt{plotblue},
\opt{plotvermillion}, \opt{plotgreen}, \opt{plotpurple}, \opt{plotsky},
\opt{plotorange}, \opt{plotyellow}, and \opt{plotgray}.

\section{Examples and Tests}

Build the documentation and example PDFs with:

\begin{lstlisting}
l3build doc
\end{lstlisting}

This creates the manual and the mode-specific example PDFs under
\filename{build/doc}. The source examples are intentionally filled documents so
title pages, headings, captions, listings, plots, tables, line numbers, links,
headers, footers, and slide furniture are visible during review.

Run the regression suite with:

\begin{lstlisting}
l3build check
\end{lstlisting}

The suite covers color mode switching, listings, plots, and smoke tests for the
CV, manuscript, thesis, and Beamer public entry points.

\section{Release Maintenance}

Before tagging a release, run the checklist in \filename{RELEASE.md}. The
short version is:

\begin{lstlisting}
l3build check
l3build doc
rsync -a build/doc/oritone.pdf doc/oritone.pdf
\end{lstlisting}

Screenshots in \filename{assets/screenshots} should be refreshed from the
current generated example PDFs whenever the visual design changes.

\end{document}
